\chapter{Walks and reachability}
\label{chap:walks}

A \emph{temporal walk} is a sequence of edges whose times are
non-decreasing (a \emph{non-strict} walk: several edges may be crossed at the
same time); a \emph{non-empty} temporal walk additionally requires at
least one edge. A walk is \emph{strict} when its times increase strictly.
Walks are formalized in \texttt{TemporalGraph/TemporalWalk.lean} and
\texttt{TemporalGraph/NonemptyTemporalWalk.lean} (which also define strictness).
Reachability, strict reachability and temporal connectivity live in
\texttt{TemporalGraph/Connectivity/Basic.lean}.

\begin{definition}[Temporal walk]
  \label{def:temporal-walk}
  \uses{def:temporal-graph}
  \lean{TemporalGraph.TemporalWalk}
  \lean{TemporalGraph.NonemptyTemporalWalk}
  \leanok
  A \emph{temporal walk} from $x$ at time $n$ to $y$ at time $m$ is either
  the empty walk (when $x = y$, $n = m$) or a non-empty temporal walk: a
  sequence of edges $(x_0,x_1,t_1), (x_1,x_2,t_2), \dots$ with
  $x_0 = x$, times non-decreasing, ending at $y$ at time $m$.
\end{definition}

\begin{definition}[Length, support, arcs of a walk]
  \label{def:walk-invariants}
  \uses{def:temporal-walk}
  \lean{TemporalGraph.TemporalWalk.length}
  \lean{TemporalGraph.TemporalWalk.support}
  \lean{TemporalGraph.TemporalWalk.arcs}
  \lean{TemporalGraph.NonemptyTemporalWalk.length}
  \lean{TemporalGraph.NonemptyTemporalWalk.support}
  \lean{TemporalGraph.NonemptyTemporalWalk.arcs}
  \leanok
  The \emph{length} of a temporal walk is its number of edges; its
  \emph{support} is the list of vertices it visits (stored from the end vertex back to the start); its
  \emph{arcs} is the list of directed edges $(x_i, x_{i+1}, t_{i+1})$ it
  traverses (stored last arc first).
\end{definition}

\begin{lemma}[Non-decreasing time]
  \label{lem:nondecreasing-time}
  \uses{def:temporal-walk}
  \lean{TemporalGraph.TemporalWalk.nondecreasing_time}
  \lean{TemporalGraph.NonemptyTemporalWalk.nondecreasing_time}
  \leanok
  A temporal walk from $x$ at time $n$ to $y$ at time $m$ satisfies
  $n \le m$.
\end{lemma}

\begin{proof}
  \leanok
  Induction on the walk: a one-edge walk starts and ends at the same time, and each further edge has a time at least that of the previous end.
\end{proof}

\begin{definition}[Reachability, temporal connectivity]
  \label{def:reachable}
  \uses{def:temporal-walk}
  \lean{TemporalGraph.Reachable}
  \lean{TemporalGraph.TemporallyConnected}
  \lean{TemporalGraph.TemporallyConnectedGraph}
  \lean{TemporalGraph.Subgraph.TemporallyConnectedSubgraph}
  \lean{TemporalGraph.Reachable.mono}
  \leanok
  $y$ is \emph{reachable} from $x$ in $G$ if there is a temporal walk from $x$ to
  $y$ (at some times $n \le m$); the walk may be empty, so every vertex reaches
  itself. A set $S$ of vertices is
  \emph{temporally connected} in $G$ (denoted $G.\mathrm{TemporallyConnected}\,S$)
  if every vertex of $S$ reaches every vertex of $S$. $G$ \emph{is temporally
  connected} (shorthand $G.\mathrm{TemporallyConnectedGraph}$) if all of $V$
  is temporally connected; a sub-temporal-graph $H$ is temporally connected
  (shorthand $H.\mathrm{TemporallyConnectedSubgraph}$) if all of $\mathrm{Verts}(H)$
  is. Reachability is \emph{monotone}: if $H \le H'$, then
  $H.\mathrm{Reachable}\,x\,y \Rightarrow H'.\mathrm{Reachable}\,x\,y$.
\end{definition}

\begin{definition}[Strict walk, strict reachability]
  \label{def:strict-reachable}
  \uses{def:temporal-walk,def:reachable}
  \lean{TemporalGraph.NonemptyTemporalWalk.IsStrict}
  \lean{TemporalGraph.TemporalWalk.IsStrict}
  \lean{TemporalGraph.StrictReachable}
  \lean{TemporalGraph.StrictReachable.mono}
  \leanok
  A temporal walk is \emph{strict} if its times increase strictly from each arc
  to the next; the empty walk is strict. $y$ is \emph{strictly reachable} from
  $x$ if some strict temporal walk goes from $x$ to $y$. Strict reachability
  implies reachability and is monotone along $\le$.
\end{definition}

\begin{lemma}[Strict and non-strict reachability agree in proper symmetric graphs]
  \label{lem:reachable-iff-strict}
  \uses{def:strict-reachable,def:is-symm-loopless,def:is-proper-happy}
  \lean{TemporalGraph.NonemptyTemporalWalk.exists_isStrict}
  \lean{TemporalGraph.reachable_iff_strictReachable}
  \leanok
  If $G$ is symmetric and proper, then $y$ is reachable from $x$ if and only if
  it is strictly reachable from $x$.
\end{lemma}

\begin{proof}
  \leanok
  Walks themselves differ: crossing an edge and immediately crossing it back at
  the same time is a non-strict walk that is not strict. Consider two consecutive
  arcs $w \to y \to z$ of a walk with the same time $t$. By symmetry, $y$ has the
  two edges $yw$ and $yz$ at time $t$, so properness forces $w = z$: the two arcs
  are a round trip, which can be deleted. The arc before them arrives at $w$ no
  later than $t$, and the arc after them leaves $z = w$ no earlier than $t$, so the
  shortened walk is still a temporal walk. Repeating this yields a strict walk,
  possibly empty. Symmetry is needed: in the oriented, proper graph with the two
  arcs $a \to y$ and $y \to z$ at time $1$, $a$ reaches $z$ but not strictly.
\end{proof}

\begin{lemma}[Cliques are temporally connected]
  \label{lem:clique-temporally-connected}
  \uses{def:reachable,def:is-simple-clique}
  \lean{TemporalGraph.temporallyConnected_of_isClique}
  \leanok
  If $S$ forms a clique in $G$ (Definition~\ref{def:is-simple-clique}), then $S$ is
  temporally connected in $G$.
\end{lemma}

\begin{proof}
  \leanok
  \uses{def:reachable}
  Two distinct vertices of $S$ are joined by an arc, which is a one-edge walk; a vertex reaches itself by the empty walk.
\end{proof}

\begin{lemma}[Reachability through an earliest edge]
  \label{lem:reachable-of-earliestEdge}
  \uses{def:reachable,def:earliest-latest-neighbor}
  \lean{TemporalGraph.reachable_of_earliestEdge}
  \lean{TemporalGraph.reachable_of_extended_earliestEdge}
  \leanok
  Assume $G$ is symmetric. If $(x,y,n)$ is the earliest edge of $x$
  (Definition~\ref{def:earliest-latest-neighbor}), every vertex reachable
  from $x$ is also reachable from $y$: prepend the edge $y \to x$ at time
  $n$ (using symmetry) to a walk out of $x$, which stays non-decreasing
  since $n$ is a lower bound on all of $x$'s times. More generally, if some
  vertex $z$ has an edge to $y$ at a time $m \le n$, then $z$ too reaches
  everything $x$ reaches, via $z \to y \to x \to \cdots$.
\end{lemma}

\begin{proof}
  \leanok
  \uses{lem:nondecreasing-time}
  Take a walk from $x$ to $z$. If it is empty, the reversed edge $y \to x$ at time $n$ (by symmetry) is a walk to $x = z$. Otherwise prepend this edge to the walk: it is admissible because $n$ is at most the time of every edge at $x$, in particular of the first edge of the walk. For the extended version, prepend $z \to y$ at time $m \le n$ followed by $y \to x$ at time $n$.
\end{proof}

\begin{lemma}[Reachability through a latest edge]
  \label{lem:reachable-of-latestEdge}
  \uses{def:reachable,def:earliest-latest-neighbor}
  \lean{TemporalGraph.reachable_of_latestEdge}
  \lean{TemporalGraph.reachable_of_extended_latestEdge}
  \leanok
  Dually: assume $G$ is symmetric. If $(x,y,n)$ is the latest edge of $x$,
  every vertex that reaches $x$ also reaches $y$: append the edge $x \to y$
  at time $n$, an upper bound on all of $x$'s times. More generally, if $y$
  has an edge to some vertex $z$ at a time $m \ge n$, then everything that
  reaches $x$ also reaches $z$, via $\cdots \to x \to y \to z$.
\end{lemma}

\begin{proof}
  \leanok
  \uses{lem:nondecreasing-time}
  Take a walk from $z$ to $x$. If it is empty, the edge $x \to y$ at time $n$ is a walk. Otherwise append this edge to the walk: it is admissible because $n$ is at least the time of every edge at $x$, in particular of the last edge of the walk. For the extended version, append $x \to y$ at time $n$ followed by $y \to z$ at time $m \ge n$.
\end{proof}
